% ===========================================================================
%  Chapitre 9 — Systèmes de numération
%  PE 2026, composante ALGÈBRE
% ===========================================================================
\bmtchapitre{Systèmes de numération}
  {Écrire un nombre dans une base demandée et effectuer des opérations en base
   quelconque, en particulier dans les systèmes binaire et hexadécimal.}
  {Algèbre \textbullet\ environ 6 heures}

Nous écrivons les nombres en base dix parce que nous avons dix doigts. Rien de
plus. Le nombre lui-même ne connaît pas cette convention : quarante-deux reste
quarante-deux, qu'on l'écrive $42$, $101010$ ou $2A$. Ce chapitre sépare enfin
le nombre de son écriture, distinction que l'on manipule sans y penser depuis
l'école primaire.

L'affaire n'est pas seulement culturelle. Un ordinateur ne dispose que de deux
états, et il compte donc en base deux ; les informaticiens abrègent en base
seize. Vous retrouverez ces deux bases dès la première année d'études
supérieures, en électronique comme en informatique. Le chapitre est court, les
techniques sont peu nombreuses, mais elles doivent devenir automatiques.

% ---------------------------------------------------------------------------
\begin{revision}
Le chapitre précédent suffit, et l'école primaire pour le reste.

\begin{enumerate}[label=\textbf{\arabic*.}, leftmargin=*, itemsep=0.35em]
  \item Effectuer la division euclidienne de $92$ par $5$, puis du quotient
        obtenu par $5$.
  \item Calculer $2^{0}$, $2^{1}$, \dots, $2^{10}$.
  \item Que signifie l'écriture $3\,405$ en base dix, si on la développe ?
  \item Calculer $16^{2}$ et $16^{3}$.
  \item Poser et effectuer l'addition $457+368$ en détaillant les retenues.
\end{enumerate}
\end{revision}

\begin{automatismes}
Sans calculatrice, sans brouillon.

\begin{enumerate}[label=\textbf{\alph*)}, leftmargin=*, itemsep=0.25em]
  \item $2^{8}$
  \item $2^{5}+2^{3}+2^{0}$
  \item $16^{2}$
  \item Reste de $92$ par $5$
  \item Reste de $255$ par $16$
  \item $3 \times 16+11$
  \item Le plus grand entier s'écrivant avec quatre chiffres en base $2$
  \item Le plus petit entier s'écrivant avec trois chiffres en base $5$
\end{enumerate}
\end{automatismes}

\begin{corrige}[automatismes]
a) $256$ \quad b) $41$ \quad c) $256$ \quad d) $2$ \quad e) $15$ \quad
f) $59$ \quad g) $15$ \quad h) $25$.
\end{corrige}

% ===========================================================================
\section{Écriture d'un entier dans une base quelconque}

\begin{definition}[écriture en base b]
Soit $b$ un entier supérieur ou égal à $2$. Tout entier naturel $N$ non nul
s'écrit de manière unique
\[
  N = a_{k}b^{k}+a_{k-1}b^{k-1}+\cdots+a_{1}b+a_{0},
\]
où les $a_{i}$ sont des entiers vérifiant $0 \leq a_{i} < b$ et
$a_{k} \neq 0$. On note alors
$N = \overline{a_{k}a_{k-1}\dots a_{1}a_{0}}^{\,b}$, et les $a_{i}$ sont les
\textbf{chiffres} de $N$ en base $b$.
\end{definition}

\begin{remarque}
Une base $b$ demande exactement $b$ chiffres, de $0$ à $b-1$. En base cinq on
n'écrit jamais le chiffre $5$ ; en base deux, seuls $0$ et $1$ existent. Un
énoncé qui propose l'écriture $\overline{153}^{\,4}$ contient donc une faute :
le chiffre $5$ n'a pas cours en base quatre.
\end{remarque}

\begin{theoreme}[existence et unicité]
L'écriture ci-dessus existe et est unique.
\end{theoreme}

\begin{demonstration}
\emph{Existence.} On construit les chiffres par divisions euclidiennes
successives : $N = bq_{0}+a_{0}$ avec $0 \leq a_{0} < b$, puis
$q_{0} = bq_{1}+a_{1}$, et ainsi de suite. La suite des quotients est
strictement décroissante et positive : elle atteint $0$, et le procédé
s'arrête.

\smallskip
\emph{Unicité.} Supposons deux écritures. Le chiffre $a_{0}$ est le reste de
$N$ dans la division par $b$, donc il est déterminé de façon unique ; le
quotient $q_{0}$ l'est aussi, et l'on conclut de proche en proche.
\end{demonstration}

\begin{visualisation}[les divisions successives, lues de bas en haut]
\begin{center}
\begin{tikzpicture}[x=1cm, y=0.86cm, font=\footnotesize]
  \node[anchor=west] at (0,3) {$92 = 5\times18+\mathbf{2}$};
  \node[anchor=west] at (0,2) {$18 = 5\times3+\mathbf{3}$};
  \node[anchor=west] at (0,1) {$\ \,3 = 5\times0+\mathbf{3}$};
  \draw[bmtLaterite, line width=1pt, -{Stealth[length=2.6mm]}]
    (4.4,3) .. controls (5.6,3) and (5.6,1.6) .. (6.55,0.35);
  \draw[bmtLaterite, line width=1pt, -{Stealth[length=2.6mm]}]
    (4.4,2) .. controls (5.4,2) and (5.4,1) .. (6.1,0.35);
  \draw[bmtLaterite, line width=1pt, -{Stealth[length=2.6mm]}]
    (4.4,1) .. controls (5.2,1) and (5.2,0.8) .. (5.65,0.35);
  \node[bmtLaterite, anchor=north, font=\normalsize\bfseries] at (6.1,0.25)
    {$332$};
  \node[bmtGris, anchor=west] at (6.9,0) {en base cinq};
\end{tikzpicture}
\end{center}

Chaque division fournit un chiffre : le premier reste donne le chiffre des
unités, le suivant celui des « cinquaines », et ainsi de suite. On lit donc
les restes \emph{de bas en haut}, ce qui est l'erreur la plus fréquente du
chapitre. Vérification immédiate :
$3\times25+3\times5+2 = 75+15+2 = 92$.
\end{visualisation}

\begin{methode}{écrire un entier dans une base b}
Diviser par $b$, noter le reste, recommencer avec le quotient jusqu'à obtenir
un quotient nul, puis lire les restes du dernier vers le premier. Terminer par
la vérification : redévelopper l'écriture obtenue et retrouver l'entier de
départ. Cette vérification coûte trois lignes et évite la moitié des erreurs.
\end{methode}

\begin{entrainement}[écrire dans une base]
\exo[\niveau{1}] Écrire $92$ en base $5$, puis $100$ en base $7$.

\exo[\niveau{1}] Écrire $2026$ en base $9$.

\exo[\niveau{2}] Développer $\overline{3\,204}^{\,5}$ et donner sa valeur en
base dix.

\exo[\niveau{2}] Combien de chiffres l'entier $1000$ a-t-il en base $2$ ? en
base $8$ ?

\exo[\niveau{3}] Déterminer la base $b$ dans laquelle $\overline{31}^{\,b}$
vaut $25$ en base dix.
\end{entrainement}

% ===========================================================================
\section{Conversions entre bases}

\begin{propriete}[passage par la base dix]
Pour convertir de la base $b$ vers la base $b'$, on développe d'abord
l'écriture en base $b$ pour obtenir la valeur en base dix, puis on effectue
les divisions successives par $b'$.
\end{propriete}

\begin{exemple}[de la base sept à la base quatre]
Convertissons $\overline{452}^{\,7}$ en base $4$.

\emph{Vers la base dix.} $4\times49+5\times7+2 = 196+35+2 = 233$.

\emph{Vers la base quatre.}
\[
\begin{aligned}
  233 &= 4\times58+1, & 58 &= 4\times14+2, \\
  14  &= 4\times3+2,  & 3  &= 4\times0+3 .
\end{aligned}
\]
En lisant les restes de bas en haut : $\overline{452}^{\,7} =
\overline{3\,221}^{\,4}$.

\emph{Vérification.} $3\times64+2\times16+2\times4+1 = 192+32+8+1 = 233$.
\end{exemple}

\begin{remarque}
Entre deux bases dont l'une est une puissance de l'autre, le passage se fait
directement, par paquets de chiffres. C'est ce qui rend le couple binaire et
hexadécimal si commode, et nous y revenons plus loin.
\end{remarque}

\begin{entrainement}[conversions]
\exo[\niveau{1}] Convertir $\overline{1\,011\,010}^{\,2}$ en base dix.

\exo[\niveau{2}] Convertir $\overline{2\,013}^{\,4}$ en base $6$.

\exo[\niveau{2}] Convertir $\overline{777}^{\,8}$ en base $2$, puis en base
dix.

\exo[\niveau{3}] Un entier s'écrit $\overline{ab}^{\,7}$ en base sept et
$\overline{ba}^{\,9}$ en base neuf. Déterminer cet entier.
\end{entrainement}

% ===========================================================================
\section{Le système binaire}

\begin{definition}[base deux]
Le \textbf{système binaire} est la numération de base $2$. Ses deux seuls
chiffres, $0$ et $1$, s'appellent des \textbf{bits}.
\end{definition}

\begin{propriete}[reconnaissance immédiate]
Un entier écrit en binaire est pair si et seulement si son dernier bit est
$0$. Plus généralement, il est divisible par $2^{k}$ si et seulement si ses
$k$ derniers bits sont nuls.
\end{propriete}

\begin{demonstration}
Tous les termes $a_{i}2^{i}$ avec $i \geq k$ sont divisibles par $2^{k}$. La
divisibilité de $N$ par $2^{k}$ équivaut donc à celle de
$a_{k-1}2^{k-1}+\cdots+a_{0}$, entier strictement inférieur à $2^{k}$ : il est
nul si et seulement si tous ses bits le sont.
\end{demonstration}

\begin{exemple}[écrire $54$ en binaire]
Les divisions successives par $2$ donnent les restes
$0,\ 1,\ 1,\ 0,\ 1,\ 1$, lus de bas en haut :
\[
  54 = \overline{110\,110}^{\,2}.
\]
On peut aussi décomposer directement : $54 = 32+16+4+2 = 2^{5}+2^{4}+2^{2}
+2^{1}$, ce qui donne les mêmes bits. Cette seconde façon est plus rapide dès
qu'on connaît les puissances de deux par cœur — et il faut les connaître.
\end{exemple}

\begin{entrainement}[binaire]
\exo[\niveau{1}] Écrire en binaire $13$, $32$ et $100$.

\exo[\niveau{1}] Convertir $\overline{1\,111}^{\,2}$ et
$\overline{100\,000}^{\,2}$ en base dix.

\exo[\niveau{2}] Combien y a-t-il d'entiers dont l'écriture binaire comporte
exactement huit bits ?

\exo[\niveau{3}] Montrer que l'écriture binaire de $2^{n}-1$ ne comporte que
des $1$.
\end{entrainement}

% ===========================================================================
\section{Le système hexadécimal}

\begin{definition}[base seize]
Le \textbf{système hexadécimal} est la numération de base $16$. Ses seize
chiffres sont $0,1,\dots,9$ puis $A$, $B$, $C$, $D$, $E$, $F$, qui valent
respectivement $10$, $11$, $12$, $13$, $14$ et $15$.
\end{definition}

\begin{propriete}[passage direct entre binaire et hexadécimal]
Puisque $16 = 2^{4}$, chaque chiffre hexadécimal correspond exactement à un
paquet de quatre bits. On convertit donc de l'un à l'autre en groupant les
bits par quatre à partir de la droite, sans passer par la base dix.
\end{propriete}

\begin{exemple}[deux conversions sans calcul]
\[
  \overline{1101\,0110}^{\,2}
  \ \longrightarrow\ \underbrace{1101}_{D}\ \underbrace{0110}_{6}
  \ =\ \overline{D6}^{\,16},
\]
et réciproquement $\overline{2F}^{\,16} = \overline{0010\,1111}^{\,2}$. En
base dix, $\overline{D6}^{\,16} = 13\times16+6 = 214$.
\end{exemple}

\begin{avertissement}
Le groupement se fait toujours \emph{à partir de la droite}. Si le nombre de
bits n'est pas multiple de quatre, on complète le paquet de gauche par des
zéros : $\overline{101\,1010}^{\,2}$ se lit $0101\ 1010$, soit
$\overline{5A}^{\,16}$.
\end{avertissement}

\begin{entrainement}[hexadécimal]
\exo[\niveau{1}] Convertir $\overline{FF}^{\,16}$ et $\overline{100}^{\,16}$
en base dix.

\exo[\niveau{1}] Écrire $255$ et $4\,096$ en hexadécimal.

\exo[\niveau{2}] Convertir $\overline{1\,0110\,1101}^{\,2}$ en hexadécimal.

\exo[\niveau{3}] Combien d'entiers s'écrivent avec exactement trois chiffres
hexadécimaux ?
\end{entrainement}

% ===========================================================================
\section{Opérations posées en base non décimale}

\begin{propriete}[addition et multiplication en base b]
Les algorithmes posés de l'école primaire restent valables dans toute base :
seule change la valeur de la retenue, qui apparaît dès que la somme partielle
atteint $b$.
\end{propriete}

\begin{exemple}[une addition posée en base cinq]
Calculons $\overline{342}^{\,5}+\overline{134}^{\,5}$.

Unités : $2+4 = 6 = 5+1$, on écrit $1$ et l'on retient $1$.
Cinquaines : $4+3+1 = 8 = 5+3$, on écrit $3$ et l'on retient $1$.
Vingt-cinquaines : $3+1+1 = 5 = 5+0$, on écrit $0$ et l'on retient $1$.
Enfin la retenue donne le chiffre de tête.
\[
  \overline{342}^{\,5}+\overline{134}^{\,5} = \overline{1\,031}^{\,5}.
\]
\emph{Vérification en base dix} : $97+44 = 141$, et
$1\times125+0+3\times5+1 = 141$.
\end{exemple}

\begin{methode}{contrôler une opération en base b}
Toujours vérifier en base dix. Convertir les deux opérandes, effectuer
l'opération en base dix, convertir le résultat et comparer. C'est plus long,
mais c'est le seul contrôle fiable tant que la base n'est pas familière.
\end{methode}

\begin{entrainement}[opérations]
\exo[\niveau{1}] Poser et effectuer
$\overline{1\,011}^{\,2}+\overline{110}^{\,2}$.

\exo[\niveau{2}] Poser et effectuer
$\overline{423}^{\,5}+\overline{344}^{\,5}$, puis vérifier en base dix.

\exo[\niveau{2}] Poser et effectuer $\overline{23}^{\,4}\times
\overline{3}^{\,4}$.

\exo[\niveau{3}] En base $b$, l'égalité $\overline{24}^{\,b}+\overline{35}^{\,b}
= \overline{61}^{\,b}$ est-elle possible ? Pour quelles valeurs de $b$ ?
\end{entrainement}

% ===========================================================================
\begin{annales}
La numération apparaît rarement seule : elle vient presque toujours greffée
sur un exercice d'arithmétique, en fin d'énoncé. Les extraits qui suivent sont
repris de sujets réellement posés ; leur provenance est indiquée à chaque
fois.

\exoann{Baccalauréat, Côte d'Ivoire, série C, session 2023 — exercice 5, question 4}
Un même entier naturel $p$ s'écrit $\overline{ac2}^{\,3}$ en base trois et
$\overline{baa}^{\,4}$ en base quatre.
\begin{enumerate}[label=\textbf{\alph*)}, leftmargin=*, itemsep=0.15em]
  \item Justifier que $3c+2$ est un multiple de $4$.
  \item En déduire que $c = 2$.
  \item Déterminer $a$ et $b$.
  \item Exprimer $p$ en base dix.
\end{enumerate}

\exoann{Baccalauréat, Madagascar, série C, session 2012 — exercice, question 2C}
Écrire l'entier $92$ dans le système de base cinq, puis vérifier le résultat
en revenant à la base dix.

\exoann{Baccalauréat, Madagascar, série C, session 2011 — exercice, partie B}
\begin{enumerate}[label=\textbf{\alph*)}, leftmargin=*, itemsep=0.15em]
  \item Convertir dans le système binaire l'entier naturel $b = 54$ écrit en
        base dix.
  \item Dresser les tables d'addition et de multiplication de $\Zn 5$.
\end{enumerate}

\exoann{Baccalauréat, Madagascar, série C, session 2026 — exercice 1}
Un entier de trois chiffres s'écrit $\overline{abc}$ en base dix. Sachant que
cet entier est divisible par $27$, écrire en base trois le plus petit d'entre
eux qui soit strictement supérieur à $100$.

\exoann{Baccalauréat, Congo-Brazzaville, série C, session 2022 — exercice 1, adapté}
Les diviseurs positifs de $195$ ayant été déterminés, écrire $195$ en base
$5$, puis en base $2$, et vérifier les deux écritures.

\exoann{Baccalauréat, Bénin, série C, session 2023 — problème 1, question 3, adapté}
Le nombre de graines semées est $p^{2}$, où $p$ est un nombre premier, et il
est compris entre $3500$ et $4000$. Après avoir déterminé ce nombre, l'écrire
en base $8$.
\end{annales}

\begin{corrige}[annales]
\textbf{1.} En base trois, $p = 9a+3c+2$ avec $a \in \{1\,;2\}$ et
$c \in \{0\,;1\,;2\}$. En base quatre, $p = 16b+4a+a = 16b+5a$ avec
$b \in \{1\,;2\,;3\}$. L'égalité donne $9a+3c+2 = 16b+5a$, soit
$3c+2 = 16b-4a = 4(4b-a)$ : c'est bien un multiple de $4$. Parmi
$c = 0,1,2$, seuls $3c+2 = 8$ convient, d'où $c = 2$. Il reste
$4a+8 = 16b$, soit $a = 4b-2$ ; la seule valeur acceptable en base trois est
$a = 2$, obtenue pour $b = 1$. Enfin $p = 9\times2+3\times2+2 = 26$, ce que
confirme $16\times1+5\times2 = 26$.

\textbf{2.} Les divisions successives donnent $92 = 5\times18+2$,
$18 = 5\times3+3$, $3 = 5\times0+3$ : en lisant les restes de bas en haut,
$92 = \overline{332}^{\,5}$. La vérification est immédiate :
$3\times25+3\times5+2 = 75+15+2 = 92$.

\textbf{3.} $54 = 32+16+4+2$, donc $54 = \overline{110\,110}^{\,2}$. Les
tables de $\Zn 5$ se remplissent en réduisant chaque somme et chaque produit
modulo $5$ ; on constate que toute classe non nulle y possède un inverse,
puisque $5$ est premier.

\textbf{4.} Le premier multiple de $27$ strictement supérieur à $100$ est
$108$. Les divisions par $3$ donnent $108 = 3\times36$, $36 = 3\times12$,
$12 = 3\times4$, $4 = 3\times1+1$, $1 = 3\times0+1$ : donc
$108 = \overline{11\,000}^{\,3}$. Vérification : $81+27 = 108$.

\textbf{5.} $195 = 5\times39+0$, $39 = 5\times7+4$, $7 = 5\times1+2$,
$1 = 5\times0+1$, d'où $195 = \overline{1\,240}^{\,5}$. En binaire,
$195 = 128+64+2+1 = \overline{1100\,0011}^{\,2}$. Les deux vérifications
redonnent bien $195$.

\textbf{6.} Le seul nombre premier compris entre $59$ et $63$ est $61$, donc
le nombre de graines vaut $3721$. En base huit : $3721 = 8\times465+1$,
$465 = 8\times58+1$, $58 = 8\times7+2$, $7 = 8\times0+7$, soit
$3721 = \overline{7\,211}^{\,8}$. Vérification :
$7\times512+2\times64+1\times8+1 = 3584+128+8+1 = 3721$.
\end{corrige}

% ===========================================================================
\begin{jecherche}[le compteur d'une machine à trier]
Une machine à trier le café compte les sacs traités et affiche ce compteur sur
huit voyants, allumés ou éteints, qui représentent une écriture binaire — le
voyant de gauche portant la plus forte valeur.

\begin{enumerate}[label=\textbf{\arabic*.}, leftmargin=*, itemsep=0.3em]
  \item Quel est le plus grand nombre de sacs que ce compteur peut afficher ?
  \item Un matin, les voyants affichent $1001\,0110$. Combien de sacs ont été
        traités ?
  \item Combien de sacs supplémentaires faut-il pour que tous les voyants
        soient allumés ?
  \item Le technicien note les relevés en hexadécimal, sur deux caractères.
        Quel relevé correspond à l'affichage de la question 2 ?
  \item Justifier que deux caractères hexadécimaux suffisent toujours, quel
        que soit l'affichage.
  \item L'entreprise veut compter jusqu'à mille sacs. Combien de voyants
        faut-il au minimum ?
\end{enumerate}
\end{jecherche}

\begin{corrige}[le compteur d'une machine à trier]
\textbf{1.} Huit bits tous allumés donnent
$2^{8}-1 = 255$ sacs.

\textbf{2.} $1001\,0110$ vaut $128+16+4+2 = 150$ sacs.

\textbf{3.} Il en faut $255-150 = 105$ de plus.

\textbf{4.} En groupant par quatre : $1001 = 9$ et $0110 = 6$, soit le relevé
$\overline{96}^{\,16}$.

\textbf{5.} Deux chiffres hexadécimaux couvrent les entiers de $0$ à
$16^{2}-1 = 255$, c'est-à-dire exactement l'étendue d'un affichage à huit
bits, puisque chaque chiffre hexadécimal absorbe quatre bits.

\textbf{6.} Il faut $2^{n}-1 \geq 1000$, soit $2^{n} \geq 1001$. Comme
$2^{9} = 512$ et $2^{10} = 1024$, dix voyants suffisent et neuf ne suffisent
pas.
\end{corrige}

% ===========================================================================
\begin{jemeteste}
Répondre par vrai ou faux, en justifiant en une ligne.

\begin{enumerate}[label=\textbf{\arabic*.}, leftmargin=*, itemsep=0.28em]
  \item L'écriture $\overline{153}^{\,4}$ a un sens.
  \item En base $b$, un entier est divisible par $b$ si et seulement si son
        dernier chiffre est $0$.
  \item $\overline{100}^{\,2}$ et $\overline{100}^{\,16}$ désignent le même
        entier.
  \item Plus la base est grande, plus l'écriture d'un entier donné est courte.
  \item Un entier pair s'écrit toujours avec un nombre pair de bits.
  \item $\overline{FF}^{\,16} = 255$.
\end{enumerate}
\end{jemeteste}

\begin{corrige}[je me teste]
\textbf{1.} Faux — en base quatre les chiffres vont de $0$ à $3$.
\textbf{2.} Vrai — tous les autres termes sont multiples de $b$.
\textbf{3.} Faux — l'un vaut $4$, l'autre $256$.
\textbf{4.} Vrai — le nombre de chiffres décroît quand la base croît, à entier
fixé.
\textbf{5.} Faux — $6 = \overline{110}^{\,2}$ s'écrit avec trois bits.
\textbf{6.} Vrai — $15\times16+15 = 255$.
\end{corrige}

% ===========================================================================
\begin{commeaubac}
\bmtsource{D'après le baccalauréat, série C, session 2023 — exercice de numération}
Un entier naturel $p$ vérifie les deux conditions suivantes : il s'écrit
$\overline{ac2}^{\,3}$ en base trois et $\overline{baa}^{\,4}$ en base quatre.

\begin{enumerate}[label=\textbf{\arabic*.}, leftmargin=*, itemsep=0.25em]
  \item Écrire les deux conditions sous forme d'égalités entre entiers, en
        précisant les valeurs possibles de $a$, $b$ et $c$.
        \hfill\textup{(0,75 pt)}
  \item Justifier que $3c+2$ est un multiple de $4$, puis déterminer $c$.
        \hfill\textup{(0,75 pt)}
  \item Déterminer $a$ et $b$, puis la valeur de $p$ en base dix.
        \hfill\textup{(1 pt)}
  \item Écrire enfin $p$ en base deux et en base seize.
        \hfill\textup{(0,5 pt)}
\end{enumerate}
\end{commeaubac}

\begin{corrige}[comme au baccalauréat]
\textbf{1.} $p = 9a+3c+2$ avec $a \in \{1\,;2\}$, $c \in \{0\,;1\,;2\}$, et
$p = 16b+5a$ avec $b \in \{1\,;2\,;3\}$, $a \leq 3$.

\textbf{2.} L'égalité $9a+3c+2 = 16b+5a$ donne $3c+2 = 4(4b-a)$. Parmi les
trois valeurs possibles de $c$, seule $c = 2$ rend $3c+2 = 8$ multiple de
$4$.

\textbf{3.} Il reste $4a+8 = 16b$, soit $a = 4b-2$. La seule valeur compatible
avec la base trois est $a = 2$, pour $b = 1$. Donc
$p = 18+6+2 = 26$.

\textbf{4.} $26 = 16+8+2 = \overline{11\,010}^{\,2}$, et en groupant par
quatre à partir de la droite, $0001\ 1010 = \overline{1A}^{\,16}$.
\end{corrige}

% ===========================================================================
\begin{concours}
\bmtsource{D'après un concours d'entrée en informatique}
Soit $b$ un entier supérieur ou égal à $2$.

\begin{enumerate}[label=\textbf{\arabic*.}, leftmargin=*, itemsep=0.25em]
  \item Montrer qu'un entier $N$ est divisible par $b-1$ si et seulement si la
        somme de ses chiffres en base $b$ l'est.
  \item Retrouver, comme cas particulier, le critère de divisibilité par $9$
        en base dix.
  \item Énoncer et démontrer le critère analogue de divisibilité par $b+1$.
\end{enumerate}
\end{concours}

\begin{corrige}[vers le concours]
\textbf{1.} Modulo $b-1$ on a $\congru{b}{1}{b-1}$, donc
$\congru{b^{i}}{1}{b-1}$ pour tout $i$. En écrivant
$N = \sum a_{i}b^{i}$, il vient $\congru{N}{\sum a_{i}}{b-1}$ : les deux
entiers ont même reste, donc $N$ est divisible par $b-1$ si et seulement si la
somme de ses chiffres l'est.

\textbf{2.} Pour $b = 10$, on obtient $b-1 = 9$ : un entier est divisible par
$9$ si et seulement si la somme de ses chiffres l'est.

\textbf{3.} Modulo $b+1$ on a $\congru{b}{-1}{b+1}$, donc
$\congru{b^{i}}{(-1)^{i}}{b+1}$. Ainsi $N$ est congru à la somme
\emph{alternée} de ses chiffres, $a_{0}-a_{1}+a_{2}-\cdots$, et $N$ est
divisible par $b+1$ si et seulement si cette somme alternée l'est. Pour
$b = 10$, c'est le critère de divisibilité par $11$.
\end{corrige}

% ===========================================================================
\begin{notepourlaclasse}
Six heures, c'est court, mais suffisant si l'on renonce à la tentation de
faire du chapitre une leçon d'informatique. Le programme demande deux choses :
écrire dans une base, et poser une opération. Le reste est du confort.

L'erreur universelle est la lecture des restes dans le mauvais sens. Imposez
dès la première séance la vérification par redéveloppement : trois lignes, et
l'erreur ne survit pas. Les élèves qui s'en dispensent sont exactement ceux
qui se trompent.

Le passage direct binaire-hexadécimal par paquets de quatre bits plaît
beaucoup et se retient sans effort. Profitez-en pour faire sentir que le choix
de la base est une convention d'écriture, jamais une propriété du nombre : la
parité, la divisibilité, la primalité ne changent pas quand on change de base.
C'est le seul message conceptuel du chapitre, et il vaut la peine d'être dit
explicitement.
\end{notepourlaclasse}
